\scanpage{22}
\begin{center}6 - 43 -\end{center}
\hfill n° 392

e) Pour que $C^*$ soit pseudo-inversible, il faut et il suffit que pour tout idéal premier $\uline{p}$ de hauteur 1 de $A$, on ait $(\uline{D}_{C/A})_{\uline{p}}:(\uline{D}_{C/A})_{\uline{p}}=C_{\uline{p}}$.

4)

a) Soient $C$ un anneau noethérien intègre, $\uline{q}$ un idéal primaire tel que $\uline{q}:\uline{q}=C$, $\uline{p}$ l'idéal premier correspondant ; on a alors $\uline{q}:\uline{p}\subset C$. Montrer que pour que l'idéal primaire $\uline{q}$ soit irréductible, il faut et il suffit que $(\uline{q}:\uline{p})/\uline{q}$ soit un $C$-module simple. (Se ramener au cas où $C$ est un anneau artinien primaire avec $\uline{q}=0$, $\uline{p}$ étant donc l'unique idéal premier de $C$ ; montrer que $0:\uline{p}$ est le socle de $C$, en notant que c'est un $(A/\uline{p})$-module, et pour que 0 soit irréductible, il faut et il suffit que le socle de $C$ soit simple).

b) On suppose que les anneaux $A$, $B$, $C$ vérifient les hypothèses de l'exerc. 2. Montrer que, pour que l'idéal fractionnaire $C^*$ de $C$ soit pseudo-inversible, il faut et il suffit que pour tout idéal premier $\uline{p}$ de hauteur 1 dans $A$, les idéaux primaires de la décomposition primaire réduite, (unique), de $\uline{p}C_{\uline{p}}$ dans $C_{\uline{p}}$, soient irréductibles. (Se ramener au cas où $A$ est un anneau de valuation discrète d'uniformisante $\pi$ ; si $\uline{m}_i$, $(1\leq i\leq r)$, sont les idéaux maximaux de $C$, pour que $C^*$ soit inversible, il faut et il suffit que les $C$-modules $C^*/\uline{m}_iC^*$ soient simples, (cf. exerc. 1 b) ou, ce qui revient au même, que les $C$-modules $(\pi C)^*/\uline{m}_i(\pi C)^*$ soient simples ; utilisant la relation $(\uline{ab})^*=\uline{a}^*:\uline{b}$, en déduire que cette condition équivaut à dire que $(\pi C:\uline{m}_i)/\pi C$ est simple pour tout $i$ ; enfin, si $\uline{q}_i$ est la composante $\uline{m}_i$-primaire de $C$, observer que $\pi C:\uline{m}_i=\uline{q}_i:\uline{m}_i$ et utiliser a)).

c) Afin que, pour tout idéal premier $\uline{p}$ de hauteur 1 de $A$, les $C_{\uline{p}}$-modules $C^*_{\uline{p}}/C_{\uline{p}}$ et $C_{\uline{p}}/(\uline{D}_{C/A})_{\uline{p}}$ soient isomorphes, il faut et il suffit que l'idéal fractionnaire $C^*$ soit pseudo-inversible. (Se ramener au cas où $A$ est un anneau de valuation discrète ; observer que si $\uline{m}_i$, $(1\leq i\leq r)$, sont les idéaux maximaux de $C$, l'hypothèse que $C^*/C$ est isomorphe à $C/\uline{D}_{C/A}$ entraîne que $C^*/\uline{m}_iC^*$ est un $C$-module simple pour tout $i$).

d) Montrer que pour que l'on ait $C_{\uline{p}}=B_{\uline{p}}$ pour tout idéal premier $\uline{p}$ de hauteur 1 de $A$, il faut et il suffit que tout idéal primaire $\neq0$ de $C$ soit pseudo-inversible. (Pour tout idéal premier $\uline{p}\neq0$ de $C$, considérer l'idéal $\uline{p}^2$).

5) Soit $K$ une extension algébrique finie de $Q$, $B$ la fermeture intégrale de $Z$ dans $K$. Afin que, pour tout sous-anneau $C$ de $B$ ayant même corps des fractions que $B$, $C^*$ soit un idéal fractionnaire inversible de $C$, il faut et il suffit que $[K:Q]\leq2$. (Si $[K:Q]=2$, montrer, en utilisant la prop.\quad, que $\uline{D}_{C/Z}$
\mspage{21}
